% Folder 57, batch 14 (archivists' pages 261-280).
% Pass: Opus 5.5 (claude-opus-5-5), 2026-10-03 — first pass, unchecked against the pages by a human.
% Skipped: 262, 264, 266, 268, 270, 272 (typed pages of a Bourbaki draft n° 227 on divisors
% and linear systems, pp. 8-12 and 44, no hand of his), 277 (blank).
% 273 is a cover sheet carrying only his title; it becomes the section heading.
% The manuscript of 261-271 runs on the versos of the typescript; 274-280 are on white paper.
% Connective prose is fast and largely illegible; formulas carry.
\documentclass[11pt,a4paper]{article}
\input{../preamble/grothendieck}

\folder{57}
\batch{14}
\pages{261}{280}
\foldertitle{Picard : tapuscrit annoté (s.d.), lettres (s.d., 1962).}
\dating{1962-[vers 1968]}
\watermark{Édition de démonstration}

\begin{document}

\page{261}
\note{le raisonnement commence avant ce lot : les numéros 1.7 à 1.11 renvoient à un § 1 qui n'est pas dans ces pages}

Démonstration algébrique \uncertain{du} \ill{} \uncertain{celle} qui sont dans le cas \uncertain{algébrique} \uncertain{sans} \uncertain{outils} \uncertain{de} RR \uncertain{analytique} …

\emph{Remarque} 1.10. \uncertain{Bien} \ill{} \ill{}, \ill{} les conditions 1.8 \add{(ou $=$ 1.9.)} \ill{} \ill{} \ill{} \ill{} \ill{} $f$ \ill{} \uncertain{projectif} (\uncertain{problèmes}). Il \ill{} \ill{} \uncertain{signale} \uncertain{ainsi} \ill{} \ill{} \ill{} \uncertain{suivant}.

$\exists$ \uncertain{partie} discrète $T$ de $Y$ telle que
\begin{enumerate}
  \item[a)] $X|(Y-T) \to Y-T$ \struck{\uncertain{fini}} \uncertain{soit} un isom.\ \ill{} ;
  \item[b)] les \uncertain{anneaux} locaux $\mathcal{O}_{Y,y}$ $(y \in T)$ soient factoriels.
\end{enumerate}

En effet, \ill{} \ill{} \ill{} \ill{} \ill{} \ill{} \uncertain{facile} de construire un \ill{} \ill{}, \ill{} \struck{\ill{}} \uncertain{analytique} \ill{} \ill{} $Y$, \emph{\uncertain{dans} \ill{} $< T$}, \uncertain{tel} que $X$, \uncertain{soit} $Y$-\uncertain{isomorphe} \ill{} \ill{} \uncertain{l'origine} \ill{} \uncertain{à} « \uncertain{éclatement} » \ill{} $Y$ \ill{} \ill{} \uncertain{générique} \ill{} « éclater » $\Sigma$.

\emph{Remarque} 1.11. Sur les variétés analytiques-formelles et algébriques formelles du \ill{} \uncertain{envisagé}. Je \ill{} \uncertain{fournir} \uncertain{cette} \uncertain{dernière} \uncertain{une} démonstration \uncertain{purement} algébrique, utilisant $\underline{\mathrm{Pic}}$, \struck{\ill{}} \add{et} \ill{} de Néron.

\page{263}
\section*{§ 2. Application à la géométrie algébrique ou analytique.}
\note{« §2 » est encadré ; le titre est de sa main}

\emph{Th.} 2.1. Soient $A$ une algèbre analytique sur $\mathbf{C}$, \add{$Y = \operatorname{Spec} A$,} $X$ un schéma propre sur $Y$, \struck{\ill{}} $\hat{Y} = \operatorname{Spf} \hat{A}$, $\hat{X} = X \times_Y \hat{Y}$, \struck{\ill{}} de sorte que [cf.\ \ill{}]
\[
\mathrm{Pic}(\hat{X}) \simeq \varprojlim \mathrm{Pic}(X_n).
\]
Alors
\begin{enumerate}
  \item[(i)] $(\mathrm{Pic}(X_n))_n$ satisfait ML ;
  \item[(ii)] $\mathrm{Pic}(X)$ et $\mathrm{Pic}(\hat{X})$ ont \uncertain{même} image dans $\mathrm{Pic}(X_n)$ [\uncertain{savoir} les images universelles] ;
  \item[(iii)] \uncertain{Pour que}
  \[
  \varphi : \mathrm{Pic}(X) \to \mathrm{Pic}(\hat{X})
  \]
  soit un isom., il suffit \uncertain{que} $H^1(X, \mathcal{O}_X)$ soit de longueur finie sur $A$, i.e.\ \uncertain{que} $\operatorname{Supp} R^1 f_*(\mathcal{O}_X) \subset \{y\}$ ($y$ le pt fermé de $Y$).
\end{enumerate}
\marginal{\uncertain{de} \uncertain{génération} \struck{\uncertain{relations}} \uncertain{ou} \uncertain{formule} \uncertain{des} \uncertain{cas}}
\note{la note marginale est reliée par une flèche à (i)}

\struck{\ill{} \ill{} \ill{} \ill{} \ill{} \ill{} \uncertain{projectif} ($\mathrm{Pic}(X_n)$) \uncertain{soit} \uncertain{cst.} \ill{}}
\note{une ligne et demie barrée de traits obliques et encadrée}

Notons \uncertain{de} suite le corollaire suivant \struck{\uncertain{conséquence}} de (i) et (ii)

\emph{Corollaire} 2.2. Conditions équivalentes

\page{265}
\begin{enumerate}
  \item[(i)] $X$ est projectif sur $Y$
  \item[(ii)] $\hat{X}$ est projectif sur $\hat{Y}$, i.e.\ $X \otimes_A \operatorname{Spec}(\hat{A})$ projectif sur $\operatorname{Spec} \hat{A}$
  \item[(iii)] \uncertain{Pour} tout entier $n$, $X_n$ est projectif sur $Y_n$.
\end{enumerate}

Évidemment (i) $\Rightarrow$ (ii) $\Rightarrow$ (iii). Prouvons (iii) $\Rightarrow$ (i). Soit $n_0$ tel que l'image de $\mathrm{Pic}(X_{n_0}) \to \mathrm{Pic}(X_0)$ soit universelle. Soit $L_{n_0}$ \uncertain{Module} \uncertain{inversible} sur $X_{n_0}$ \add{\uncertain{ample}} \marginal{$L_0 = L_{n_0}|X_0$} relativement $/Y_{n_0}$. \add{Donc $\mathcal{E}$ \uncertain{se} \uncertain{remonte} \uncertain{et} 2.1.\ (ii)} \struck{Je \uncertain{prends} \ill{} d'un} \ill{} \uncertain{module} \add{\uncertain{inversible}} $\mathcal{E}$ sur $\hat{X}$, \uncertain{donc} \uncertain{en} \uncertain{vertu} \uncertain{de} 2.1.\ (ii) il \uncertain{provient} \uncertain{aussi} d'un \uncertain{module} \uncertain{inversible} \add{\uncertain{ample}} $L$ sur $X$. \uncertain{Le} \uncertain{dernier} \uncertain{est} \uncertain{ample} \uncertain{en} \uncertain{vertu} de EGA III.\ \uncertain{Cqfd}.

Notons maintenant que 2.1.\ \struck{\ill{}} se déduit aisément de 1.7.\ par la GAGA de Serre-Grauert-Remmert-Grothendieck.

\page{267}
\emph{Corollaire} 2.3. Supposons $\dim X_0 \leqslant 1$. Alors $X$ est projectif sur $Y$.
\note{chacun des corollaires 2.3 et 2.4 est marqué d'un trait vertical à gauche}

\emph{Corollaire} 2.4. Soit $A$ \uncertain{une} \uncertain{algèbre} \uncertain{analytique} \add{et \emph{réduite},} \uncertain{normale} \add{\uncertain{soit} $Y = \operatorname{Spec} A$,} \add{\uncertain{supposons} $Y$ à singularité \emph{isolée}} \add{\uncertain{soient}} $\tilde{Y} = \operatorname{Spec}(\hat{A})$, $y$ ($\tilde{y}$) le pt fermé de $A$. Alors l'hom.
\[
\mathrm{Pic}(Y_y) \longrightarrow \mathrm{Pic}(\tilde{Y}_{\tilde{y}}) \quad \text{est un \add{\emph{iso}}\struck{\emph{monomorphisme}}morphisme.}
\]
\note{les variantes « et réduit », « Supposons $Y$ à singularité isolée » sont ajoutées au-dessus de la ligne et encadrées ; l'ordre de lecture des insertions n'est pas sûr}

\struck{Cet hom.\ est injectif}

On peut supposer que \uncertain{les} \uncertain{singularités} \uncertain{isolées} de $Y$ \uncertain{soient} de dim.\ $\geqslant 2$. Alors, \uncertain{toute} \uncertain{classe} \uncertain{de} \uncertain{module} \uncertain{inversible} $E'$ sur $Y' = Y - \{y\}$, $i_*(E')$ est coh.\ (\uncertain{où} $i : Y' \to Y$ \uncertain{l'immersion} can.). En particulier, $i_*(\mathcal{O}_{Y'})$ est cohérent, \uncertain{donc} \uncertain{une} algèbre $A'$ finie sur $A$. \struck{\uncertain{Remarquons} que \ill{}} Cette dernière \uncertain{remarque} \uncertain{s'applique} \uncertain{aussi} \uncertain{à} $Y_1$ fini sur \struck{$Y$}, réduit, tel que $Y'_1 = Y_1|Y' \to Y'$ soit un isom., [\uncertain{donc} $\tilde{Y}'_1 = Y_1|\tilde{Y}' \to \tilde{Y}'$ est un isom.], \struck{\ill{} \ill{} $\Gamma($} et \struck{\ill{}} \ill{}
\[
\Gamma(Y_1) \to \Gamma(Y'_1) \quad \text{est un \uncertain{isom.}}
\]
\uncertain{Cela} \uncertain{nous} \uncertain{ramène} \uncertain{au} cas où $\Gamma(Y_1) \to \Gamma(Y')$ est un \uncertain{isom.}, i.e.\ où \uncertain{prof} $A \geqslant 2$. Un raisonnement facile
\marginal{\uncertain{déjà} \uncertain{vu} \ill{} \ill{} \ill{} \ill{} \uncertain{cas} \uncertain{où} $A$ \uncertain{est} \uncertain{non} \uncertain{normal} \ill{}}
\drawing{dans la marge gauche, une figure en éventail : quatre triangles adjacents partant d'un même sommet, coupés par un arc}

\page{269}
(\uncertain{comme} \uncertain{plus} \uncertain{haut} \uncertain{dans} \uncertain{le} \uncertain{cas} \uncertain{précédent}) \uncertain{montre} \uncertain{que}
\[
\mathrm{Pic}(Y') \to \mathrm{Pic}(\tilde{Y}') \quad\text{est \emph{injectif}.}
\]
Reste à prouver la surjectivité. Pour ceci, \uncertain{on} \uncertain{utilise} \uncertain{la} \emph{résolution des singularités} (Hironaka) : il existe $X \to Y$ \uncertain{morphisme} \uncertain{propre} \uncertain{dominant}, $X$ \struck{\ill{}} \uncertain{non} \uncertain{singulier}, tel que $X' = X|Y' \xrightarrow{\ \sim\ } Y'$, et $X = \overline{X'}$. Donc on a
\[
\mathrm{Pic}(\text{\struck{$X$}}\,Y') \simeq \mathrm{Pic}(X') \simeq \mathrm{Pic}(X)/I \simeq \mathbf{Z}^{I}
\]
où $I$ est l'ensemble \uncertain{des} \uncertain{comp.} \uncertain{irréd.} de $X - X' = X_0$. \uncertain{On} \uncertain{a} \uncertain{aussi} \uncertain{de} \uncertain{même} [\uncertain{notant} \uncertain{que} $\tilde{X} = X \times_Y \tilde{Y}$ est \uncertain{régulier}, \uncertain{en} \uncertain{plus} \uncertain{de} \struck{\ill{}} \uncertain{est} \uncertain{régulier} \uncertain{sur} $X$ \uncertain{puisque} \uncertain{est} \uncertain{régulier}]
\[
\mathrm{Pic}(\tilde{Y}') \simeq \mathrm{Pic}(\tilde{X}') \simeq \mathrm{Pic}(\tilde{X})/I \simeq \mathbf{Z}^{I}
\]
\uncertain{On} \uncertain{conclut} \uncertain{par} \uncertain{le} th.\ 2.1.\ (iii) \uncertain{qui} \uncertain{donne} $\mathrm{Pic}(\tilde{X}) \simeq \mathrm{Pic}(X)$, d'où la conclusion voulue.
\note{l'écriture « $\mathbf{Z}^{I}$ » suit la page ; dans ces deux formules le quotient « $/I$ » est à entendre comme quotient par le sous-groupe engendré par les composantes}

\page{271}
\emph{Questions.} \struck{\ill{}} Soit $A$ un \uncertain{anneau} local \add{\uncertain{excellent}} noethérien. \uncertain{On} \uncertain{suppose} $A$ \emph{\uncertain{hensélien}}, \uncertain{et} \uncertain{un} « \uncertain{bon} \uncertain{anneau} », i.e.\ $\operatorname{Spec}\hat{A} \to \operatorname{Spec} A$ est \struck{\ill{}} \uncertain{universellement} \uncertain{régulier}. Soit $X$ propre sur $A$, et $\hat{X} = X \otimes_A \hat{A}$.
\begin{enumerate}
  \item[(i)] Est-il vrai que les images universelles dans $\mathrm{Pic}(X_n)$ [\uncertain{dans} \uncertain{un} \uncertain{sens} \struck{\ill{}} \uncertain{évident}] \uncertain{de} $\operatorname{Im}(\mathrm{Pic}(\hat{X}) \to \mathrm{Pic}(X_n))$ est \uncertain{dans} $\operatorname{Im}(\mathrm{Pic}(X) \to \mathrm{Pic}(X_n))$ ? En particulier, si $\hat{X}/\hat{Y}$ est projectif, est-il vrai que $X/Y$ est projectif ? \struck{\ill{}} \add{p.\ ex.} \uncertain{l'anneau} \uncertain{hensélisé} \uncertain{d'un} \uncertain{anneau} local de la géom.\ alg.\ sur $\mathbf{C}$ ?
  \item[(ii)] Supposons que $\operatorname{Supp} R^1 f_*(\mathcal{O}_X) \subset \{y\}$, \struck{\ill{}} p.\ ex.\ $X|(Y-\{y\}) \xrightarrow{\ \sim\ } Y - \{y\}$. Est-il vrai alors que
  \[
  \mathrm{Pic}(X) \xrightarrow{\ \sim\ } \mathrm{Pic}(\hat{X}) \ ??
  \]
  \item[(iii)] \struck{N.B.} Supposons que $Y$ soit à singularité isolée en $y$. Est-il vrai que
  \[
  \mathrm{Pic}(Y - \{y\}) \xrightarrow{\ \sim\ } \mathrm{Pic}(\hat{Y} - \{\hat{y}\})
  \]
\end{enumerate}
[Si (ii) est vrai, et si on peut désingulariser, alors (iii) est vrai …]. Comment \uncertain{démontrer} \uncertain{ici} \uncertain{à} \uncertain{moins} \add{\ill{}} \uncertain{que} \uncertain{l'idéal} \uncertain{de} \uncertain{sing} $(Y)$, \uncertain{dans} \uncertain{le} \uncertain{cas} \uncertain{des} \uncertain{anneaux} alg.\ \uncertain{analytiques} ?
\marginal{What does Murre \uncertain{do} in « Publications » p.~16 \uncertain{\textit{iiiitis}} ?? \uncertain{chiar}}
\note{la note marginale, en anglais, est écrite en oblique dans la marge gauche et rattachée par une accolade aux questions (ii) et (iii) ; « Publications » renvoie sans doute aux Publications mathématiques de l'IHÉS}
\note{une grande boucle relie « universellement régulier » et la fin de (i) : la phrase sur l'anneau hensélisé semble être un exemple ajouté pour la question (i)}

\section{Schémas de Néron-Severi. Polarisation}
\note{titre de sa main, écrit seul en haut d'un feuillet par ailleurs blanc (p. 273) : « Schémas de Néron-Severi », et au-dessus « Polarisation ». Les pages qui suivent sont sur un autre papier, blanc, et d'une autre encre que les pages 261 à 271}

\page{274}
\note{dans ces pages, les foncteurs et préschémas soulignés par lui sont rendus par $\underline{\mathrm{Pic}}$, $\underline{\mathrm{N\acute{e}r}}$ (un trait) et $\underline{\underline{\mathrm{Pic}}}$, $\underline{\underline{\mathrm{N\acute{e}r}}}$ (deux traits), comme sur la page ; l'exposant noté ici ${}^{*}$ est un petit signe dont la forme n'est pas sûre}
\emph{Définition.} $X$ préschéma sur $S$. Un élément $\xi$ de $\mathrm{Pic}(X/S)$ est dit \emph{numériquement équivalent à $0$} si sa restriction à toute fibre $X_s$ est numériquement équivalente à zéro [i.e. un multiple d'icelle peut se déformer dans $0$]. Deux éléments $\xi, \xi'$ de $\mathrm{Pic}(X/S)$ sont dits numériquement équivalents si $\xi - \xi'$ est numériquement équivalent à $0$.

Soit $\mathrm{Pic}^{\tau}(X/S) \subset \mathrm{Pic}^{*}(X/S)$ \uncertain{l'ensemble} des $\xi$ numériquement équivalents à $0$, et
\[
\mathrm{N\acute{e}r}'(X/S) = \mathrm{Pic}(X/S)/\mathrm{Pic}^{\tau}(X/S).
\]
Notons : si \struck{$S' \to S$ est fid.\ plat [et quasi-compact] et} $S' \to S$ est un changement de base, d'où $\mathrm{Pic}(X/S) \to \mathrm{Pic}(X'/S')$, il induit $\mathrm{Pic}^{\tau}(X/S) \to \mathrm{Pic}^{\tau}(X'/S')$, et si $S' \to S$ est surjectif, alors \struck{un élément $\xi$ de} $\mathrm{Pic}^{\tau}(X/S)$ est exactement l'image inverse de $\mathrm{Pic}^{\tau}(X'/S')$ par $\mathrm{Pic}^{*}(X/S) \to \mathrm{Pic}(X'/S')$, d'où \uncertain{dans} ce cas
\[
\mathrm{N\acute{e}r}'(X/S) \longrightarrow \mathrm{N\acute{e}r}'(X'/S') \quad\text{est \emph{injectif}.}
\]

\note{le carré qui suit est dessiné dans la marge gauche, en face du changement de base}

\begin{tikzcd}
X \arrow[d] & X' \arrow[l] \arrow[d] \\
S & S' \arrow[l]
\end{tikzcd}

Considérons $S' \mapsto \mathrm{N\acute{e}r}'(X'/S')$ comme foncteur contravariant
\[
(\mathrm{Sch})^{\circ}_{/S} \longrightarrow (\mathrm{Ens}),
\]
et rendons-le \emph{exact pour les morphismes fid.\ plats} \add{loc.\ de type} \struck{quasi-compacts} \emph{fini} [donc pour la localisation habituelle aussi]. On trouve un foncteur $S' \mapsto \mathrm{N\acute{e}r}(X'/S')$, et en particulier un \add{ensemble} $\mathrm{N\acute{e}r}(X/S)$, \uncertain{ensemble} des polarisations de $X$ sur $S$.

\page{275}
\emph{Interprétation} lorsque $\underline{\mathrm{Pic}}_{X/S}$ existe, ($\underline{\mathrm{Pic}}^{\tau}_{X/S}$ ouvert !). Alors la donnée d'un élément de $\mathrm{Pic}''(X/S)$ \struck{\ill{}} équivaut à la donnée d'un sous-préschéma \add{ouvert} $P$ de $\underline{\mathrm{Pic}}^{\tau}_{X/S}$ qui soit un fibré principal homogène \emph{trivial} sous $\underline{\mathrm{Pic}}^{\tau}_{X/S}$. Alors $\mathrm{N\acute{e}r}(X/S)$ est l'ensemble des sous-préschémas \struck{$\longrightarrow$} de $\underline{\mathrm{Pic}}_{X/S}$ qui sont formellement principaux homogènes $P$ \struck{sous} sous $\underline{\mathrm{Pic}}^{\tau}_{X/S}$, et \ill{} par \ill{} \uncertain{catégorie} \uncertain{loc.} de type fini. [N.B. \struck{Lorsque} Ces sous-préschémas sont nécessairement \emph{ouverts}. Si $\underline{\mathrm{Pic}}^{\tau}_{X/S}$ est plat et de type fini sur $S$, \uncertain{alors} la condition \ill{} \ill{} \uncertain{sous} $P$ qu'on vient de mettre signifie aussi que $P$ est \emph{fid.\ plat} et de type fini sur $S$]. Par suite,
\marginal{de type fini ?}

\struck{dire que lorsque}

\emph{Proposition.} \emph{Supposons que $\underline{\mathrm{Pic}}_{X/S}$ existe, et que $\underline{\mathrm{Pic}}^{\tau}_{X/S}$ soit fid.\ plat [de type fini] sur $S$.} Pour que le foncteur $\mathrm{N\acute{e}r}_{X/S}$ \struck{\ill{}} dans $(\mathrm{Sch})^{\circ}_{/S}$ soit représentable, il f.\ et s.\ que le quotient
\[
\underline{\mathrm{N\acute{e}r}}_{X/S} = \underline{\mathrm{Pic}}^{*}_{X/S}/\underline{\mathrm{Pic}}^{\tau}_{X/S}
\]
existe, et \struck{alors} que $\underline{\mathrm{Pic}}_{X/S} \to \underline{\mathrm{N\acute{e}r}}_{X/S}$ soit \emph{fid.\ plat} \struck{et quasi-compact} \add{et de type fini} de noyau $\underline{\mathrm{Pic}}^{\tau}_{X/S}$,
\note{la proposition est encadrée d'un trait vertical à gauche, et la phrase continue en haut de la page suivante}

\page{276}
et alors $\underline{\underline{\mathrm{N\acute{e}r}}}_{X/S}$ est représenté par $\underline{\mathrm{N\acute{e}r}}_{X/S}$.

\emph{Corollaire 1.} S'il en est ainsi, $\underline{\mathrm{N\acute{e}r}}_{X/S}$ est \struck{représenté} séparé ssi $\underline{\mathrm{Pic}}^{\tau}_{X/S}$ est fermé dans $\underline{\mathrm{Pic}}_{X/S}$. En tout cas, $\underline{\mathrm{N\acute{e}r}}_{X/S}$ est \emph{non ramifié} sur $S$.

\emph{Corollaire 2.} Supposons que $\underline{\mathrm{Pic}}^{\tau}_{X/S}$ soit \emph{propre} \add{et \emph{plat}} sur $S$, et que $\underline{\mathrm{Pic}}_{X/S}$ soit réunion \add{croissante} d'ouverts \add{$U_i$} stables par $\underline{\mathrm{Pic}}^{\tau}_{X/S}$ et quasi-projectifs sur $S$. [Cette dernière condition est conséquence de $\underline{\mathrm{Pic}}^{\tau}_{X/S}$ de type fini sur $S$, dans le cas où $X$ est projectif, plat, \add{\struck{\ill{}} \ill{} \uncertain{universellement}} (et \uncertain{primaire} sur $S$]. Alors $\underline{\mathrm{N\acute{e}r}}_{X/S}$ existe et est séparé non ramifié sur $S$.
\marginal{La \uncertain{première} condition \ill{} \ill{} \uncertain{vérifiée} \ill{} de $X/S$ \ill{} \uncertain{sans torsion} \ill{} des $\mathrm{Pic}^0$ propres…}
\note{les deux corollaires sont marqués chacun d'un trait vertical à gauche}

Il semble que dès que $\underline{\mathrm{Pic}}^{\tau}_{X/S}$ est \emph{plat} \struck{\ill{}} \add{et fermé dans $\underline{\mathrm{Pic}}_{X/S}$} sur $S$, $\underline{\mathrm{N\acute{e}r}}_{X/S}$ doive exister : c'est lié à la question de la caractérisation des foncteurs représentables par un objet non ramifié sur $S$ … . Au commencement ! [À vrai dire, même l'existence de $\underline{\mathrm{Pic}}^{\tau}_{X/S}$ ne semble pas indispensable !]
\marginal{?}
\note{l'insertion « et fermé dans $\underline{\mathrm{Pic}}_{X/S}$ » est entourée et rattachée par un trait ; le paragraphe est marqué en marge d'un trait vertical et d'un point d'interrogation}

\page{278}
\subsection{Schémas polarisés}

Soit $X$ projectif sur $S$ loc.\ noeth., \add{plat et \uncertain{Steinien},} \struck{\ill{}} \struck{polarisation stricte}
\marginal{semble excessif, par exemple pour l'énoncé de K.\ et Matsusaka}
\note{la note marginale est reliée par un trait au mot entouré « Steinien »}

\struck{\ill{} $(X/S)$ \ill{} $X$ est le domaine d'un élément de $\mathrm{Pic}(X)$, \ill{} qui induit un diagramme \ill{} d'un faisceau inversible \ill{} \ill{}}
\note{deux lignes encadrées et barrées de traits obliques}

On suppose que $\underline{\underline{\mathrm{Pic}}}_{X/S}$ existe. Soit $\mathrm{N\acute{e}r}'(X/S)$ l'ensemble quotient de $\mathrm{Pic}(X)^{\tau/S}$ \struck{\ill{}} par $\mathrm{Pic}^{\tau}(X)^{\tau/S}$, où $\mathrm{Pic}(X)^{\tau/S}$ est l'image inverse de $\Gamma(S, \underline{\underline{\mathrm{Pic}}}^{\tau}_{X/S})$ par $\mathrm{Pic}(X) \to \Gamma(S, \underline{\underline{\mathrm{Pic}}}_{X/S})$.
\marginal{\struck{\ill{} \ill{} \ill{} \ill{} \ill{} $X/S$}}
\note{les exposants de $\mathrm{Pic}(X)$ sont lus sous réserve : le premier est en partie raturé}

Les éléments de $\mathrm{N\acute{e}r}'(X/S)$ s'appellent \struck{pseudo-} \struck{polarisations} \add{\uncertain{néronisations}} \emph{strictes}. Les \uncertain{vraies} : c'est un foncteur contravariant en $T$, soit $\mathrm{N\acute{e}r}'_{X/S}$, \struck{multiplicatif} \uncertain{transformant} \struck{\ill{}} les sommes en produits. Rendons ce foncteur compatible avec les descentes fid.\ plates de type fini ; on trouve un foncteur $\underline{\mathrm{N\acute{e}r}}''_{X/S}$. \struck{Un élément de la \ill{}}
\marginal{[i.e. \uncertain{les morphismes} fid.\ plats de type fini $T' \to T$ \uncertain{sont tels que} $\underline{\mathrm{N\acute{e}r}}'_{X/S}(T) \to \underline{\mathrm{N\acute{e}r}}'_{X/S}(T')$ soit injectif ; \uncertain{i.e.} les morphismes fid.\ plats de type fini sont de descente \uncertain{stricte} pour $\underline{\mathrm{Pic}}'_{X/S}$]}

\[
\underline{\mathrm{N\acute{e}r}}''_{X/S}(S) = \mathrm{N\acute{e}r}''(X/S),
\]
les éléments de cet ensemble s'appellent les \struck{pseudo-} \struck{prépolarisations} \add{\uncertain{néronisations}} \emph{de $X$ sur $S$}. Donc :
\[
\text{(\ill{})}\quad \underline{\mathrm{N\acute{e}r}}''_{X/S}(T) \simeq \mathrm{N\acute{e}r}''(X_T/T) \quad \text{(isom.\ can.)},
\]
une \uncertain{polarisation} de $X/S$ est définie par la donnée d'une polarisation stricte \add{i.e.\ d'un $\underline{L}$ inversible} $\xi$ d'un $X_T/T$, \struck{\ill{}} sur $X_T/T$, \add{où $T \to S$ fid.\ plat de type fini,} tel que les deux images inverses de $\xi$ sur $X_{T \times_S T}/T \times_S T$ soient identiques, i.e.\ tel que $p_1^{*}(\underline{L})\, p_2^{*}(\underline{L})^{-1}$ soit de \uncertain{torsion} \add{\uncertain{équivalente à $0$}} \uncertain{sur} $T \times_S T$.

\page{279}
$\xi$ de $\mathrm{N\acute{e}r}'(X_T/T)$ et $\xi'$ de $\mathrm{N\acute{e}r}'(X_{T'}/T')$ définissent la même \uncertain{pseudopolarisation} \uncertain{ssi} leurs images inverses sur $\mathrm{N\acute{e}r}'(X_{T \times_S T'}/T \times_S T')$ sont égales, i.e.\ si $p_1(\underline{L})\, p_2(\underline{L}')^{-1}$ est \uncertain{numériquement équivalent} à $0$ \uncertain{rel.} à $T \times_S T'$.

Si p.\ ex.\ $\underline{\underline{\mathrm{Pic}}}^{\tau}_{X/S}$ est plat sur $S$ et si le quotient $\underline{\underline{\mathrm{N\acute{e}r}}}_{X/S} = \underline{\underline{\mathrm{Pic}}}_{X/S}/\underline{\underline{\mathrm{Pic}}}^{\tau}_{X/S}$ existe « raisonnablement » [i.e.\ \struck{\ill{}} \uncertain{que} \uncertain{la relation} d'équivalence \uncertain{définie} par $\underline{\underline{\mathrm{Pic}}}^{\tau}_{X/S}$ \uncertain{soit} effective, \ill{} la projection $\underline{\underline{\mathrm{Pic}}}_{X/S} \to \underline{\underline{\mathrm{Pic}}}_{X/S}/\underline{\underline{\mathrm{Pic}}}^{\tau}_{X/S}$ \uncertain{fid.\ plate}] — p.\ ex.\ si $\underline{\underline{\mathrm{Pic}}}^{\tau}_{X/S}$ est propre sur $S$ — alors \struck{\ill{}} $\underline{\mathrm{N\acute{e}r}}''_{X/S}$ s'identifie au foncteur \uncertain{associé} à $\underline{\underline{\mathrm{N\acute{e}r}}}_{X/S}$, les \uncertain{néronisations} de $X/S$ sont les sections de $\underline{\underline{\mathrm{N\acute{e}r}}}_{X/S}$ sur $S$, \struck{\ill{}} \ill{} justifie la \uncertain{terminologie} \struck{\ill{}} \add{\ill{} \ill{}}.
\struck{Deux}

\struck{Deux néronisations \ill{} \ill{} \ill{} \uncertain{équivalentes}}

\struck{Si} On définit un sous-ensemble $\mathrm{N\acute{e}r}'^{+}(X/S)$ de $\mathrm{N\acute{e}r}'(X/S)$ par les \uncertain{classes} des faisceaux inversibles sur $X$ \uncertain{relativement} amples \uncertain{rel.} à $S$, d'où un sous-foncteur $\underline{\mathrm{N\acute{e}r}}'^{+}_{X/S}$, donnant naissance à un sous-foncteur $\underline{\mathrm{N\acute{e}r}}^{+\prime\prime}_{X/S}$ de $\underline{\mathrm{N\acute{e}r}}''_{X/S}$, qui \uncertain{représente} \uncertain{explicitement} \ill{} les « \ill{} \ill{} » \ill{} « faisceaux \ill{} inversibles amples ». Dans le cas particulier

\page{280}
envisagé où le schéma de Néron $\underline{\underline{\mathrm{N\acute{e}r}}}_{X/S}$ existe, le foncteur $\underline{\mathrm{N\acute{e}r}}^{+\prime\prime}_{X/S}$ est associé au sous-\uncertain{schéma} $\underline{\underline{\mathrm{N\acute{e}r}}}^{+}_{X/S}$ de $\underline{\underline{\mathrm{N\acute{e}r}}}_{X/S}$, \uncertain{image} de $\underline{\underline{\mathrm{Pic}}}^{+}_{X/S}$. Un élément \ill{} est appelé une \emph{polarisation} \struck{\ill{}} de $X/S$. Elles forment un \struck{pseudo-}monoïde (\uncertain{sans élément neutre}) commutatif. \struck{Deux} \add{Deux} polarisations $\xi, \xi'$ sont dites \emph{équivalentes} s'il existe des entiers $n, n' > 0$ tels que $n\xi = n'\xi'$. Les \add{classes de} polarisations correspondent aux \uncertain{demi-droites} d'un cône convexe $C$ \struck{[\uncertain{ouvert}]} dans un espace vectoriel $V$ sur $\mathbf{Q}$, \uncertain{engendré} par $C$, \uncertain{donc} les \uncertain{points} de $C$ \uncertain{sont} « \uncertain{rationnels} » \ill{}.
\marginal{\struck{\ill{} \ill{} \ill{}} \ill{} $\mathrm{N\acute{e}r}''(X/S)$).}

Polarisation \emph{très ample} (\uncertain{pour} une polarisation donnée, il en existe \uncertain{toujours} \uncertain{une} \uncertain{équivalente} qui est \emph{très ample}).

\note{la page 280 s'arrête ici ; l'exposé sur les polarisations continue sans doute au-delà de ce lot}

\end{document}
